\section{Autoregressive image modeling: por qué una interfaz unificada no ganó en todas las tareas}

Si una image puede discretizarse en tokens, la idea más natural es seguir el camino de GPT: text-to-image, image-to-text e image completion son solo distintos ordenamientos de tokens y distintas mask. Esta línea tiene la elegancia de una interfaz unificada, pero la high-resolution sequence length, la tokenizer information loss y la sequential decoding la ponen frente a una fuerte competencia de diffusion en calidad de generación y en velocidad.

\subsection{CogView/DALL-E/Parti: primero convertir la imagen en una secuencia discreta}

Sea un image tokenizer que codifica la imagen $I$ en codes $z=(z_1,\ldots,z_m)$, y sean $x$ los tokens del texto. La factorization de text-to-image es
\[
p(z\mid x)=\prod_{i=1}^{m}p(z_i\mid x,z_{<i}).
\]
Para entrenar basta con concatenar los text tokens y los image tokens en una sola sequence, pero el número de image tokens puede llegar a miles, y la generación debe recorrer en serie todas las posiciones.

\lecturefigure{slide-22-autoregressive-text-to-image.jpg}{Autoregressive text-to-image: image tokenizer más GPT sequence}{Grabación oficial de Stanford Online 00:46:47}

\begin{importantbox}{El image tokenizer aporta a la vez compresión y pérdida}
El VQ tokenizer comprime los píxeles continuos en un codebook finito, lo que acorta mucho la secuencia; pero el texto fino, las texturas, la geometría y la precisión del color pueden perderse en la etapa de codificación. El Transformer posterior no puede recuperar información que el tokenizer nunca conservó.
\end{importantbox}

\teachervoice{00:45:00--00:49:32, el ponente coloca CogView y DALL-E/Parti en una misma recipe sencilla: primero discretizar y luego autoregress. La sencillez es una ventaja, y también el origen de los cuellos de botella en velocidad y en detalle.}

\subsection{Universal modeling: poder hacer todas las direcciones no equivale a ser el mejor en cada una}

Un mismo sequence model puede expresar varias tareas mediante ordenamientos y mask:
\[
\begin{aligned}
\text{text}\rightarrow\text{image}:&\ [x;z],\\
\text{image}\rightarrow\text{text}:&\ [z;x],\\
\text{masked image}:&\ [z_{\text{visible}};z_{\text{target}}].
\end{aligned}
\]
Esto sí unifica la API, pero en image understanding puede quedar por debajo de un VLM que conserva features continuas, y en image generation puede quedar por debajo de diffusion.

\lecturefigure{slide-23-universal-modeling.jpg}{Universal vision-language modeling: unificación por ordenamiento y compromiso de desempeño}{Grabación oficial de Stanford Online 00:49:32}

\begin{warningbox}{Un objective unificado y una solución óptima unificada son dos cosas distintas}
Que un modelo pueda expresar la likelihood de varias tareas no significa que, con parámetros, datos y compute limitados, alcance a la vez la Pareto frontier de cada uno de los specialized systems. Una representación compartida aporta transfer, pero también interference y bottleneck inadecuados.
\end{warningbox}

\teachervoice{00:49:32--00:52:11, el ponente hace una valoración muy mesurada de la línea de universal modeling en la que él mismo participó: logró la unificación, pero en image generation no supera a diffusion, y en image understanding tampoco supera a los VLM con características visuales continuas.}

\subsection{Resumen de la sección}

La mayor ventaja de autoregressive image modeling es la interfaz unificada y la madura language-model machinery; sus mayores problemas son la pérdida del tokenizer discreto, la generación en serie de secuencias largas y la linealización de las relaciones espaciales bidimensionales. No es que no pueda generar buenas imágenes, sino que la practical frontier se desplazó hacia diffusion, que se adapta mejor al cómputo paralelo de imágenes.
